% 《大学物理简明教程》第5版第11章狭义相对论计算11.3.10—11.3.18
% 印刷/PDF第291页。原题数值逐页核对，研究性独立推导。
\subsection{解答题：相对论动能、质能关系与核反应}

\sourceproblem{11.3.10}{291}{291}
\problemtext 某粒子的动能是其静止能量的$n$倍，求运动速度。
\workings
$K=(\gamma-1)m_0c^2=nm_0c^2$，
从而$\gamma=n+1$，
\[
\boxed{v=c\sqrt{1-\frac1{(n+1)^2}}}.
\]
\crosscheck $n=0$时$v=0$；$n\to\infty$时渐近$c$，始终不超光速。

\sourceproblem{11.3.11}{291}{291}
\problemtext 电子速率$v=0.99c$，其静止能量约$0.51$ MeV。
求相对论动能。
\workings
$\gamma=(1-0.99^2)^{-1/2}\approx7.0888$，
\[
\boxed{K=(\gamma-1)m_ec^2
=(7.0888-1)(0.51)
\approx3.105\,\mathrm{MeV}}.
\]
\crosscheck 不可使用$\frac12m_ev^2$的非相对论公式，
因为此处$v$已接近$c$。

\sourceproblem{11.3.12}{291}{291}
\problemtext 电子分别(1)从静止加速到$0.1c$，
(2)从静止加速到$0.9c$，求相应做功。
\workings
由相对论功—能定理，
\[
W(v)=\Delta K=(\gamma(v)-1)m_ec^2,
\quad m_ec^2\simeq0.51\,\mathrm{MeV}.
\]
在$\beta=0.1$时
\[
\boxed{W_{0.1c}=
\left(\frac1{\sqrt{0.99}}-1\right)0.51
\approx2.57\times10^{-3}\,\mathrm{MeV}
=2.57\,\mathrm{keV}}.
\]
在$\beta=0.9$时
\[
\boxed{W_{0.9c}=
\left(\frac1{\sqrt{0.19}}-1\right)0.51
\approx0.660\,\mathrm{MeV}}.
\]
\crosscheck 两个结果都用\emph{静止到末速}的能量增量；
第二问不能用$0.9c$的经典动能公式。

\sourceproblem{11.3.13}{291}{291}
\problemtext 物体的相对论质量相对静质量增大10\%，
问沿运动方向的长度比静止长度缩短多少百分比？
\workings
按原书“运动质量”定义$\gamma m_0=1.10m_0$，
即$\gamma=1.10$。长度收缩
\[
\frac{L}{L_0}=\frac1\gamma=\frac1{1.10}
\approx0.90909,
\]
故
\[
\boxed{\frac{L_0-L}{L_0}=1-\frac1{1.1}
\approx9.09\%}.
\]
\crosscheck 不是缩短10\%，因为长度缩短因子与质量增加因子互为倒数。

\sourceproblem{11.3.14}{291}{291}
\problemtext 电子速率由$1.2\times10^8\,\mathrm{m/s}$增至
$2.4\times10^8\,\mathrm{m/s}$，求外力所做功。
\workings
以$c=3.0\times10^8$ m/s，速度分别为$0.4c,0.8c$，
$\gamma_1=1/\sqrt{0.84}$、$\gamma_2=1/\sqrt{0.36}=5/3$。
功等于动能改变量
\[
\boxed{W=(\gamma_2-\gamma_1)m_ec^2
=\left(\frac53-\frac1{\sqrt{0.84}}\right)0.51\,\mathrm{MeV}
\approx0.294\,\mathrm{MeV}}.
\]
换算SI单位约$4.72\times10^{-14}$ J。
\crosscheck 应当相减两末态的总能量，不可把两动能直接用经典公式求差。

\sourceproblem{11.3.15}{291}{291}
\problemtext 粒子静止质量为$m_0$、动能恰等于自身静能。
求其相对论总能量、相对论质量以及动量大小。
\workings
$K=(\gamma-1)m_0c^2=m_0c^2$，故$\gamma=2$，
从而
\[
\boxed{E=2m_0c^2},\qquad
\boxed{m_{\mathrm{rel}}=2m_0},\qquad
\boxed{v=\frac{\sqrt3}{2}c}.
\]
动量$p=\gamma m_0v$，得
\[
\boxed{p=\sqrt3\,m_0c}.
\]
\crosscheck 不变量关系
$E^2-(pc)^2=4m_0^2c^4-3m_0^2c^4=m_0^2c^4$正确。
现代通常不用“运动质量”一词；此处为对原教材术语的明确定义。

\sourceproblem{11.3.16}{291}{291}
\problemtext 太阳每秒向外辐射能量约$5.0\times10^{26}$ J，
估算每秒质量减少量，并以太阳当前质量$2\times10^{30}$ kg
估算维持相同辐射功率可持续的理论时间。
\workings
由$E=\Delta mc^2$，
\[
\boxed{\dot m=\frac{P}{c^2}
=\frac{5.0\times10^{26}}{(3.0\times10^8)^2}
\approx5.56\times10^9\,\mathrm{kg/s}}.
\]
如果假设全部太阳现有质量均可转换为这种辐射且功率不变，
\[
\boxed{t_{\rm formal}=\frac{Mc^2}{P}
=\frac{2\times10^{30}\cdot9\times10^{16}}{5\times10^{26}}
=3.6\times10^{20}\,\mathrm s\approx1.14\times10^{13}\,\mathrm{yr}}.
\]
\crosscheck 这只是\emph{把全部静质量转成辐射}的数学上限估计，
并不是实际恒星核燃料寿命；真实太阳仅部分质量可参与核燃烧且亮度随演化变化。

\sourceproblem{11.3.17}{291}{291}
\problemtext 某核电站年发电量约$10^{10}$ kWh，
即$3.6\times10^{16}$ J。若电能全部由核材料静质量减少提供，
所消耗的质量是多少？
\workings
按100\%质量—电能转换效率的理想极限，
\[
\boxed{\Delta m=\frac{E}{c^2}
=\frac{3.6\times10^{16}}{9.0\times10^{16}}
=0.40\,\mathrm{kg}}.
\]
\crosscheck 实际反应燃料质量远大于0.40kg，
因为核反应只有微小质量亏损且热—电转换效率也小于100\%。

\sourceproblem{11.3.18}{291}{291}
\problemtext 热核反应$\,{}^{2}_{1}\mathrm H+{}^{3}_{1}\mathrm H
\to{}^{4}_{2}\mathrm{He}+{}^{1}_{0}\mathrm n$。
给出各粒子静质量：
$m_D=3.3437\times10^{-27}$kg、
$m_T=5.0049\times10^{-27}$kg、
$m_{\mathrm{He}}=6.6425\times10^{-27}$kg、
$m_n=1.6750\times10^{-27}$kg，
求释放能量。
\workings
反应前后静质量差为
\[
\Delta m=m_D+m_T-m_{\mathrm{He}}-m_n
=(3.3437+5.0049-6.6425-1.6750)\times10^{-27}
=3.11\times10^{-29}\,\mathrm{kg}.
\]
释放能量
\[
\boxed{Q=\Delta mc^2
=(3.11\times10^{-29})(3\times10^8)^2
\approx2.80\times10^{-12}\,\mathrm J}.
\]
按$1\,\mathrm{MeV}=1.602176634\times10^{-13}$ J，
\[
\boxed{Q\approx17.5\,\mathrm{MeV}}.
\]
\crosscheck 两个质量和的差必须保留$10^{-27}$kg的共同数量级；
最终正的质量亏损说明反应放能。
